\section{Obstructions that rule out several tempting shortcuts}
\label{sec:logical-obstructions}

These examples distinguish failures of a proposed proof transfer from
impossibility claims about all alternative calculi. They leave open a
new calibration using additional strength of a redesigned seed theory.

\subsection{The bundle's gate cannot simply be erased}

Classically the actual bundle is a conjunction of its seed component
and
\[
 R(a,b)\to Q_cR(a,b).
\]
Its native implementation writes this conjunction through implication
and negation, but it has the same truth condition. Consequently an
invalid vacuity conjunct cannot be repaired by the other conjunct,
even when the same code $c$ also names the seed matrix.

For an explicit collision, take native labels $c=a=1$ and $b=2$.
Interpret membership in any model of ZFC and assign $x_1=\varnothing$,
$x_2=\{\varnothing\}$. The antecedent $x_1\in x_2$ is true. Varying
$x_1$ to $\{\varnothing\}$ makes it false, so
$Q_1R(1,2)$ is false. The whole ungated bundle is therefore false,
regardless of its seed component; the additional fact that matrix code
$c=1$ is native falsity does not alter this conclusion. The existing
set-model soundness proof genuinely needs the gate or a different
bundle construction.

\subsection{Cofinally fresh vacuity does not give all fresh vacuity}

One proposed substitute emits $R(a,b)\to Q_cR(a,b)$ only when
$c>\max(a,b)$, or only when $c>a+b$. Neither restriction recovers the
full logical vacuity family from the retained logical rows alone.

Here is an explicit assignment semantics. Let the domain have at least
three elements. For an assignment $s:\mathbb N\to D$, define
\[
 \mu_s(v)=\min\{i:s(i)=v\}
\]
whenever $v$ is among the values read by the displayed atom, and set
\[
 R(a,b)[s]\quad\Longleftrightarrow\quad
 \mu_s(s(a))<\mu_s(s(b)).
\]
Interpret equality as actual equality, implication classically, and
$Q_i$ as quantification over replacement of coordinate $i$ of $s$.
The relation $R(a,a)$ is always false. Equality congruence holds in
both argument positions because equal coordinate values have equal
first-occurrence indices. All the retained quantifier and equality
rows, the merged equality row with a fresh binder, and classical
implicational rows are valid under this semantics.

If $c>\max(a,b)$, changing $s(c)$ cannot alter either first-occurrence
index: each is already attained at an index at most $a$ or $b$.
Thus every such restricted vacuity instance is valid. However, fix any
$m$, put $a=m+2$, $b=m+3$, and $c=m$, and use three distinct domain
values $A,B,C$. Assign all coordinates below $m$ the value $C$, and put
\[
 s(m)=A,\quad s(m+1)=B,\quad s(a)=A,\quad s(b)=B.
\]
Then $R(a,b)$ is true. Changing coordinate $m$ from $A$ to $B$ changes
the first occurrences to $\mu(A)=m+2$ and $\mu(B)=m$, making the atom
false. Hence $R(a,b)\to Q_cR(a,b)$ fails even though $c\ne a,b$.
The counterexample works above every prescribed finite prefix of labels.
It also satisfies the still weaker family with fixed native roles
$a=0,b=1$ and $c\geq2$.

This model refutes the suggested derivation from the \emph{logical}
fragment. It is not claimed to satisfy the full seed theory. A scheme
with fixed seed roles and restricted vacuity would need a new argument
using its nonlogical axioms; an appeal to unrestricted variable
renaming alone is insufficient.

\subsection{Both directions of membership congruence occur in the proof}

The current lower calibration uses both
\[
 \rho_3:\ E(a,b)\to(R(c,a)\to R(c,b)),\qquad
 \rho_5:\ E(c,a)\to(R(c,b)\to R(a,b)).
\]
They are independent over the other logical ingredients. To see why
$\rho_5$ cannot simply be dropped and later recovered, choose a family
of unary predicates $P_i$ with $P_0$ empty and $P_1$ the whole nonempty
domain, and interpret
\[
 R(i,j)[s]\quad\Longleftrightarrow\quad P_i(s(j)).
\]
Actual equality and ordinary assignment quantifiers validate $\rho_3$,
all fresh vacuity instances, and the other retained logical rows.
But $\rho_5$ fails when its equality premise identifies the values at
labels $0$ and $1$: the predicates attached to the labels remain
different. For the other direction use
$R(i,j)[s]\Longleftrightarrow P_j(s(i))$, which validates $\rho_5$ and
can falsify $\rho_3$. In either interpretation $R(0,0)$ is always false,
so the native falsity alias is respected.

There is also an exact dependency audit of the existing proof.
\texttt{KThmK.rel\_L} is $\rho_5$ and
\texttt{KThmK.rel\_R} is $\rho_3$;
\texttt{atom\_subst}, \texttt{atom\_subst\_rev}, and the atomic case of
\texttt{leibniz} use both.
In \texttt{KEind222.Eq\_Eind}, the incoming-membership component uses
$\rho_3$ and the outgoing-membership component uses $\rho_5$.
In \texttt{KIndisc222.qmem}, quotient membership uses both components,
first \texttt{hx.2 y} and then \texttt{hy.1 x'}.
Thus the current quotient and first-order completeness route does not
justify deleting either row. A different seed-theoretic quotient
argument could still be a worthwhile separate research direction.

\paragraph{Dependency pins and proof status.}
The relevant modules are
\texttt{KThmK.lean}, \texttt{KEind222.lean}, and
\texttt{KIndisc222.lean} in the repository's set-theory closure-kernel
formalization, read at commit
\texttt{2762b8a183830b1f5a30cf059ffb4d92e77cddde}; the exact paths are
listed with the binder theorem's dependencies. The new countermodels
and scheduler lemmas are ordinary mathematical arguments. They were
not newly compiled by a proof-assistant kernel in this work.
